% --- Archon physics formalization source begin ---
% archon:physics
% archon:chemistry
% archon:covers IChO2026Problems/problem_icho_2026_t3_a6.lean
% archon:source-report reports/icho_2026/problem_icho_2026_t3_a6.source.json
% archon:problem-id icho_2026_t3
% archon:part-id T3-A6

\section{Icho Chemistry problem icho\_2026\_t3\_a6}

\paragraph{Problem source.}
\#\# Source
58th International Chemistry Olympiad, Tashkent, Uzbekistan, 2026.
Official-materials index: https://www.icho2026.uz/problems
English problem PDF: ../icho\_2026\_source/raw/theory\_problem.pdf
English solution/rubric PDF: ../icho\_2026\_source/raw/theory\_solution.pdf
Problem PDF page 30; printed local question page 6; solution starts on PDF page 30.
The page PNG is primary visual evidence for formulas, structures, tables, and figures.

\#\# T3. Into Reticular Chemistry
T3. Into Reticular Chemistry 7\% Covalent organic frameworks (COFs) are a class of polymers that form two- or three-dimensional structures through reactions between organic precursors, resulting in covalent bonds to afford porous, crystalline materials. Two examples of honeycomb type 2D-COFs formed by boronate ester and boroxine rings (repeat units indicated using dashed lines) are shown in the figure below. The black and red bold lines represent different building blocks in all COF topology figures. 3.1 Determine the empirical formula of COF-1 and calculate the mass percentage (\%, to two decimal places) of carbon in COF-1. 3.2 Calculate the internal diameter, d, in Å, of the honeycomb for COF-2, with the following assumptions: • C–C/C=C (arenes) = 1.39 Å • C–B = 1.56 Å • B–O = 1.38 Å • Note, the diameter of a circle, 𝑑, inscribed inside a hexagon, is given by 𝑑= √3𝑎, where 𝑎is the side length of the hexagon. • Neglect the width of the linkers. The structures shown below can form 2D and 3D COFs. For this question consider only formation of COFs with boronate esters, imines, and C=C bonds through condensation reactions. A–D represent compound classes with different functional groups. Tetragonal means Tetragonal, Hexagonal means Hexagonal, Trigonal means Trigonal, Kagome means Kagome, Tetrahedral means Tetrahedral. 3.3 Fill all the table cells with either new examples of monomer combinations that give the topologies shown above or with “XXX” if you are certain that no additional combinations satisfy that topology. Both the combinations and the “XXX” entries will be graded. Incorrect answers (both combinations and “XXX” entries) will be penalised. Blank table cells will not be graded. Tetragonal 1 Tetragonal 2 Hexagonal 1 Hexagonal 2 Trigonal Kagome Tetrahedral A2 + B3 E1 + D2 C2 + D4 Whilst imine-based COFs are easily synthesised, they are also vulnerable to hydrolysis under acidic and basic conditions. The addition of hydroxy groups to the aldehyde or amine monomers can overcome this. IR spectrum of COF-3 contained stretches corresponding to the imine functional group, but COF-4 did not. COF-4 contains a different type of C=N bonds, and elemental analysis revealed a loss of six additional hydrogen atoms per repeat unit relative to COF-3 as the only change. COF-5 irreversibly isomerises to COF-6, which did not contain imine stretches in its IR spectrum. 3.4 Draw the repeat units for COFs 3–6. Indicate the edges of each repeat unit with a dashed line as shown for COF-1 below. Macrocycle X was synthesised by post-synthetic modification of COF-7 (C2 + D4) followed by degradation. 3.5 Draw the smallest repeat unit of macrocycle X. π-π interactions result in different stacking modes of COF layers, some of which are shown below. Depending on stacking geometry, the interaction energy can be different as shown below. 𝐸/ kJ mol−1 Interaction geometry benzene (X = CH) – benzene (b-b) −7.9 −12.6 benzene-triazine (X = N) (b-t) −49.8 −55.6 triazine–triazine (t-t) −6.7 −16.7 COF-8, synthesised from (E1 + D2), is notable for its remarkable stability. The AA' slightly shifted mode of a COF-8 bilayer has a π-π stacking interaction energy of –67.1 kJ mol−1 between two layers of one repeat unit.

\#\# Current subquestion 3.6
Calculate the π-π stacking energy between two layers of one repeat unit of COF-8, for the AA, AB and AB' arrangements. Assume π-π stacking happens only between aromatic units.

\#\# Dataset metadata
IChO 2026; T3; raw rubric points=12; kind=theory; formalization\_ready=true.

\paragraph{Shared problem context.}
T3. Into Reticular Chemistry 7\% Covalent organic frameworks (COFs) are a class of polymers that form two- or three-dimensional structures through reactions between organic precursors, resulting in covalent bonds to afford porous, crystalline materials. Two examples of honeycomb type 2D-COFs formed by boronate ester and boroxine rings (repeat units indicated using dashed lines) are shown in the figure below. The black and red bold lines represent different building blocks in all COF topology figures. 3.1 Determine the empirical formula of COF-1 and calculate the mass percentage (\%, to two decimal places) of carbon in COF-1. 3.2 Calculate the internal diameter, d, in Å, of the honeycomb for COF-2, with the following assumptions: • C–C/C=C (arenes) = 1.39 Å • C–B = 1.56 Å • B–O = 1.38 Å • Note, the diameter of a circle, 𝑑, inscribed inside a hexagon, is given by 𝑑= √3𝑎, where 𝑎is the side length of the hexagon. • Neglect the width of the linkers. The structures shown below can form 2D and 3D COFs. For this question consider only formation of COFs with boronate esters, imines, and C=C bonds through condensation reactions. A–D represent compound classes with different functional groups. Tetragonal means Tetragonal, Hexagonal means Hexagonal, Trigonal means Trigonal, Kagome means Kagome, Tetrahedral means Tetrahedral. 3.3 Fill all the table cells with either new examples of monomer combinations that give the topologies shown above or with “XXX” if you are certain that no additional combinations satisfy that topology. Both the combinations and the “XXX” entries will be graded. Incorrect answers (both combinations and “XXX” entries) will be penalised. Blank table cells will not be graded. Tetragonal 1 Tetragonal 2 Hexagonal 1 Hexagonal 2 Trigonal Kagome Tetrahedral A2 + B3 E1 + D2 C2 + D4 Whilst imine-based COFs are easily synthesised, they are also vulnerable to hydrolysis under acidic and basic conditions. The addition of hydroxy groups to the aldehyde or amine monomers can overcome this. IR spectrum of COF-3 contained stretches corresponding to the imine functional group, but COF-4 did not. COF-4 contains a different type of C=N bonds, and elemental analysis revealed a loss of six additional hydrogen atoms per repeat unit relative to COF-3 as the only change. COF-5 irreversibly isomerises to COF-6, which did not contain imine stretches in its IR spectrum. 3.4 Draw the repeat units for COFs 3–6. Indicate the edges of each repeat unit with a dashed line as shown for COF-1 below. Macrocycle X was synthesised by post-synthetic modification of COF-7 (C2 + D4) followed by degradation. 3.5 Draw the smallest repeat unit of macrocycle X. π-π interactions result in different stacking modes of COF layers, some of which are shown below. Depending on stacking geometry, the interaction energy can be different as shown below. 𝐸/ kJ mol−1 Interaction geometry benzene (X = CH) – benzene (b-b) −7.9 −12.6 benzene-triazine (X = N) (b-t) −49.8 −55.6 triazine–triazine (t-t) −6.7 −16.7 COF-8, synthesised from (E1 + D2), is notable for its remarkable stability. The AA' slightly shifted mode of a COF-8 bilayer has a π-π stacking interaction energy of –67.1 kJ mol−1 between two layers of one repeat unit.

\paragraph{Current subquestion.}
Calculate the π-π stacking energy between two layers of one repeat unit of COF-8, for the AA, AB and AB' arrangements. Assume π-π stacking happens only between aromatic units.

\paragraph{Recorded answer/context.}
COF-8 represents hexagonal 2 topology. Each repeat unit contain 4 benzene and 1 triazine rings. For AA and AA’ stacking modes there will be 4 b-b and 1 t-t interactions. For AB each repeat unit will have 1 b-t interactions as shown below for AB’ arrangement (hexagons is redrawn to show all interactions). Thus, calculations will be as follows: Interaction number Interaction energy, kJ mol−1 Total energy, kJ mol−1 b-b b-t t-t b-b b-t t-t AA 4 0 1 −7.9 −49.8 −6.7 −38.3 AA' 4 0 1 −12.6 −55.6 −16.7 –67.1 AB 0 1 0 −7.9 −49.8 −6.7 −49.8 AB' 0 1 0 −12.6 −55.6 −16.7 –55.6 𝐴𝐴∶4 ⋅(−7.9) + 1 ⋅(−6.7) = −38.3 kJ mol−1 𝐴𝐵∶1 ⋅(−49.8) = −49.8 kJ mol−1 𝐴𝐵′ ∶1 ⋅(−55.6) = −55.6 kJ mol−1 (AB') interlayer stacking mode energy matches slipped benzene-triazine π-π stacking. 4 pt for each correct total energy in AA, AB, AB' 12.0 pt

\paragraph{Figure/image paths.}
\begin{itemize}
\item ../icho\_2026\_source/image/T3\_page-6.png
\item ../icho\_2026\_source/image/T3\_page-5.png
\end{itemize}

\paragraph{Formalization target.}
create a compiling Lean file with sorry bodies at `IChO2026Problems/problem\_icho\_2026\_t3\_a6.lean`.
The Lean declarations must preserve every mathematical object, quantifier, hypothesis, side condition, approximation/error guarantee, and requested conclusion from the source statement.
Use Mathlib/Physlib/Chemistry names found through LeanExplore where available. Prefer the configured benchmark Base library; introduce a local definition only when the source genuinely requires an object that the environment does not expose.

\begin{theorem}\leanok
[Icho Chemistry formalization target]
\label{thm:physics:icho_2026_t3_a6:target}
\lean{IChO2026Problems.T3A6.cof8_AA_AB_ABPrime_stacking_energies}
\uses{def:icho_2026_t3_a6:stacking_energy}
The source interaction-energy table and stacking diagram determine the three
COF-8 stacking energies.
\end{theorem}
\begin{proof}

\leanok
Match each named AA, AB, and AB$'$ geometry to its table entry and use the
source diagram to select the observed AA$'$ interaction where required.
\end{proof}

\begin{definition}\leanok
[COF-8 stacking energy]
\label{def:icho_2026_t3_a6:stacking_energy}
\lean{IChO2026Problems.T3A6.cof8StackingEnergy}
For one stacking mode, the COF-8 interaction energy is the sum over aromatic
pair types of the contact count times the energy selected by that mode's
geometry.
\end{definition}
\begin{proof}

\leanok
Each aromatic contact contributes its tabulated pair energy once, so summing
count times energy over the three pair types gives the bilayer total.
\end{proof}
% --- Archon physics formalization source end ---
